The four integer forms have exact degrees, not merely syntactic upper
bounds. Every fixed numeral has degree zero. For the retained-$f$ family,
the relevant register degrees are
\[
\begin{array}{c|rrrrrrrrrrrr}
\text{register}&a&H&\Delta&R&k&Q&V&C_y&U&Z&S&D_\tau\\\hline
\deg&4&4&8&5&5&10&7&24&6&13&28&22.
\end{array}
\]
The weighted radicands in $N_4$ have degrees $14,34,40$, with auxiliary
numerator degree $36$ and denominator degree $10$.

Write a subscript zero for the leading homogeneous part. Useful expressions are
\begin{align*}
q_0&=6\beta j,&W_0&=12w,&Y_0&=q_0v,&X_0&=q_0W_0,\\
a_0&=X_0Y_0,&i_0&=2r,&R_0&=-q_0^3(W_0C-q_0z),&k_0&=R_0+ha_0,\\
U_0&=8q_0^2a_0,&M_0&=2q_0^2,&Q_0&=a_0^2f^2,&V_0&=-f^2R_0,\\
S_0&=-24Y_0^4k_0^4,&(D_\tau)_0&=X_0^2Y_0^4k_0^2.&&&
\end{align*}

\subsection{Four roots and three roots}
In the four-root norm the guard root has degree 20, strictly exceeding
the degrees 17, 7, 13 of the other roots and degree 13 of the center.
Its leading term is therefore
\[
 (N_4)_{400}=Q_0^8(81U_0^2S_0)^8\ne0.
\]
Multiplication by $F_f$, of degree 14 with leading part
$-a_0^2i_0^2c^4$, gives exact degree 414. A coefficient certificate is
\[
 [j^{196}w^{66}v^{98}h^{32}r^2f^{16}c^4]P_{14}
 =-2^{402}3^{302}\beta^{196}\ne0.
\]
For three roots, the first Pell root is uniquely largest, with degree 17.
Hence $N_3$ has degree $8\cdot17+4\cdot10=176$ and leading part
\[
 Q_0^4(9U_0^2(D_\tau)_0)^4.
\]
The factors $F_f,F_g$ add degrees 14 and 28. Thus $P_{15}$ has degree
218, with the explicit coefficient
\[
 [j^{104}w^{38}v^{50}r^2f^8h^{12}c^4]P_{15}
 =-2^{209}3^{151}\beta^{104}\ne0.
\]
These certificates remain nonzero on the actual fixed numerical tuple.

\subsection{The two-root equation and degree 124}
The center $Z$ contains the leading approximation to the auxiliary root.
Write
\[
 Z=UV+s,\qquad s=Un+Mb+H(v-q^2),\qquad\deg s=7.
\]
Then the exact ring identity
\begin{equation}
 \delta(Z^2-M^2D_m)-U^2C_y
 =\delta(s^2+2UVs-M^2D_m)-U^2(V^2-1)
 \label{eq:centeridentity}
\end{equation}
lowers the degree of this inner expression to 30. Therefore $N_2$ has
degree at most 60. Its complete degree-60 part is
\[
 4Q_0^2U_0^2V_0^2(s_0^2-A_0),\quad
 s_0=M_0a_0(4n+c),\quad A_0=M_0^2a_0^2c^2,
\]
equivalently
\[
 (N_2)_{60}=2^{13}q_0^8a_0^8f^8R_0^2n(2n+c)\ne0.
\]
The remaining factors have degrees 14, 28, and 22. Consequently
\[
 \deg P_{16}=60+14+28+22=124.
\]
For a single-coefficient certificate,
\[
 [j^{58}w^{18}v^{24}r^2h^6f^8z^2n^2c^4]P_{16}
 =2^{113}3^{77}\beta^{58}>0.
\]
The powers $h^6,z^2,n^2$ isolate the corresponding terms in
$k_0^6,R_0^2,n(2n+c)$ uniquely. This proves exact degree independently
of any algebraic relations among the other five fixed coefficients.

\subsection{Five roots and degree 1072}
Without retaining $f$, the degrees change to
$\deg Q=22$, $\deg V=19$, $\deg C_y=60$, $\deg Z=25$.
The auxiliary radical satisfies
\[
 C_y/\delta=V^2+\frac{V^2-1}{Q-1}.
\]
For a rigorous homogeneous-degree computation, scale all inputs by a
formal parameter and work in its Laurent-series extension. The two
auxiliary roots are $\pm(V+\text{terms of degree }-3)$.
After multiplication by $U$, their correction has degree 3.
Sixteen conjugate factors retain the degree-25 auxiliary term; in the
other sixteen it cancels the center's $UV$, leaving the guard root of
degree 20 as the unique leading term. Clearing denominators contributes
$16\cdot22$. Hence
\[
 \deg P_{13}=16(25+20+22)=1072.
\]
Its leading form is
$Q_0^{16}(2U_0V_0)^{16}(729U_0^2S_0)^8$, now with
$Q_0=a_0^4i_0^2c^4$ and $V_0=-a_0^2i_0^2c^4R_0$ from the
unretained-strong definitions. The explicit certificate is
\[
 [j^{480}w^{160}v^{192}h^{32}r^{64}c^{128}z^{16}]P_{13}
 =2^{1000}3^{696}\beta^{480}>0.
\]

\subsection{Literal arithmetic evaluation}
The complete DAG counts are 195, 148, 130, and 118 operations for
13, 14, 15, and 16 integers. They include the finalizer and all scalar
products. The product factors are evaluated after cancellation of their
constant terms, as displayed in Section~\ref{sec:integer}.
The five-root evaluator also shares terms in the first biquadratic
square when its two linear components coincide. These are exact ring
identities, not on-solution substitutions. The circuit files enumerate
every operation, and the verifier checks topology and liveness.

